\section{Data}
\label{sec:data}

\subsection{Sources and retrieval}

All statistics are season aggregates for the English Premier League, retrieved for the
\SeasonPrimary{} season, which the study is about, and the \SeasonReplication{} season,
which is used only for replication. Retrieval goes through the \texttt{soccerdata}
library, which caches every page to disk and rate limits requests to one every seven
seconds, so no request is ever issued twice and the source site is not burdened. Raw
pages are cached before any transformation, and every retrieved table is written to
disk untouched before a single column is renamed.

Two properties of the retrieval path are worth recording, because both took work to
establish and both affect what any replication of this study will encounter. The first
is that FBref sits behind a challenge that returns a refusal to ordinary HTTP clients,
so only the library release that drives a real browser can reach it. The second is that
this release, while it can reach the site, exposes only five of the eleven season stat
tables that the site publishes. We restore the missing six by subclassing the reader and
re-declaring the table identifiers, reusing the library's own downloader, cache and
parser rather than writing any request loop of our own. Establishing the correct
identifiers required inspecting live pages, which is how we found that the goal and shot
creation table is served under an identifier that does not match its documented name.

\subsection{The withdrawal of advanced statistics}
\label{sec:data:withdrawal}

On 20 January 2026 the operator of FBref removed every statistic derived from its
commercial data provider, following termination of the licence and a demand that the
data be deleted. The removal was retroactive rather than prospective, so it applies to
completed seasons as well as future ones.

The practical consequence is more subtle than a set of missing tables, and it is a trap
for anyone building on this source. The affected tables are still served, still carry
their full column headers, and still return one row per player. Only the values are
gone. A feature list assembled from the site's published schema therefore looks entirely
correct and yields a feature matrix that is empty. Table~\ref{tab:availability} reports
the audit we ran to establish exactly what survived, counting non-null values for every
column of every table in both seasons.

\begin{table}[t]
\centering
\small
\caption{Populated columns in each FBref player table. The tables are still served with complete headers, so the empty counts are the number of columns that resolve correctly and contain no values.}
\label{tab:availability}
\begin{tabular}{lrrrr}
\toprule
 & \multicolumn{2}{c}{2024-25} & \multicolumn{2}{c}{2025-26} \\
\cmidrule(lr){2-3}\cmidrule(lr){4-5}
Table & Columns & Empty & Columns & Empty \\
\midrule
standard & 17 & 0 & 17 & 0 \\
shooting & 11 & 0 & 11 & 0 \\
passing & 21 & 19 & 21 & 19 \\
passing types & 16 & 14 & 16 & 14 \\
goal shot creation & 17 & 16 & 17 & 16 \\
defense & 17 & 14 & 17 & 14 \\
possession & 21 & 20 & 21 & 20 \\
playing time & 17 & 0 & 17 & 0 \\
misc & 13 & 2 & 13 & 2 \\
keeper & 19 & 0 & 19 & 0 \\
keeper adv & 26 & 23 & 26 & 23 \\
\bottomrule
\end{tabular}
\end{table}


What remains is close to the statistical vocabulary of two decades ago: appearances,
minutes, goals, assists, shots, shots on target, cards, fouls committed and drawn,
offsides, crosses, interceptions and tackles won, together with basic goalkeeping
volume. What disappeared is the entire modern apparatus, including expected goals,
expected assisted goals, shot and goal creating actions, every touch and carry and
take-on count, all passing volume and progression, and, for goalkeepers, post-shot
expected goals, sweeper actions, distribution length and cross handling.

\subsection{Restoring shot quality}

Because a role discovery study with no measure of chance quality would be badly
impoverished, we merge a second and independent source. Understat publishes its own
expected goals model and continues to provide non-penalty expected goals, expected
assists, key passes, shot counts, and the two possession involvement measures xGChain
and xGBuildup. These restore the creation and shot quality block that the withdrawal
removed, at the cost of introducing a second provider whose model differs from the one
that produced the original FBref values. We treat that as a limitation rather than a
repair, and we never mix an Understat quantity with an FBref quantity that measures the
same thing.

Records are matched across the two sources within club, using three passes of
decreasing strictness: exact match on a normalised name, then a token subset rule that
resolves short forms against full names, then an approximate string match. The token
subset pass is accepted only when exactly one candidate matches, so a club fielding
several players who share a given name can never be silently conflated. Seven clubs are
named differently by the two sites and are mapped explicitly. The procedure matches
\MatchRatePrimary{} percent of \SeasonPrimary{} records and \MatchRateRepl{} percent of
\SeasonReplication{} records. Unmatched players retain missing values for the Understat
block and are counted rather than quietly dropped.

\subsection{Filtering, normalisation and feature construction}

The \SeasonPrimary{} pull contains \NPlayersRawPrimary{} player records, of which
\NOutfieldPrimary{} are outfield players and \NKeepersPrimary{} are goalkeepers. A
player who changed club during the season appears once per club in the source. We retain
those rows for reference but conduct every analysis on an aggregated row per player, in
which counting statistics are summed and rates are recomputed from their components
rather than averaged across spells of unequal length, since averaging a rate over
unequal exposure is not the rate. \NMultiClubPrimary{} players required this treatment
in \SeasonPrimary{}.

Players with fewer than \MinMinutes{} minutes, five full matches, are excluded from
clustering while being retained in the descriptive tables. The threshold trades sample
size against the noise in a per-90 rate computed over a short exposure, and
Appendix~\ref{sec:appendix:sensitivity} reports the same analysis at 270 and 900 minutes.
At the chosen threshold \NEligiblePrimary{} outfield players are eligible in
\SeasonPrimary{} and \NEligibleRepl{} in \SeasonReplication{}, distributed across
\NDFPrimary{} defenders, \NMFPrimary{} midfielders and \NFWPrimary{} forwards.

Counting statistics are divided by nineties played. Quantities that are already rates,
such as shot accuracy, are left alone rather than being divided a second time. The three
statistics that can only be recorded while the opponent has the ball, interceptions,
tackles won and fouls committed, are additionally corrected for how much time a player's
team spends out of possession, by the procedure set out in
Section~\ref{sec:methods:padj}. Minutes
and nineties are retained as context columns and never used as features, since squad
usage describes a player's opportunity rather than the role he performs within it. For
the same reason the team success columns, points per match and goal difference while on
the pitch, are excluded from every feature set: they measure the quality of the squad a
player happens to belong to.

Positions are parsed from the source's comma separated string. The first listed position
gives a primary position and a position group in defender, midfielder, forward and
goalkeeper, while the full string is retained for the analysis of players listed at more
than one position. The final outfield feature set contains \NFeatures{} features and is
listed in full in Appendix~\ref{sec:appendix:features}. Every feature in it was confirmed
to hold non-null values in both seasons before inclusion, which is the only reliable
defence against the empty column problem described in
Section~\ref{sec:data:withdrawal}. Features are standardised twice, once across the whole
outfield pool for the global analysis and once within position group for the
position stratified analysis, and both versions are retained.

\subsection{Known limitations of the data}

Season aggregates discard context. A tackle count does not record where on the pitch the
tackles happened, an expected goals total does not distinguish a striker who takes six
low quality shots from one who takes two good ones, and the possession correction described in
Section~\ref{sec:methods:padj} reaches only the three statistics recorded while the
opponent has the ball, leaving the attacking statistics still sensitive to how much of it
a player's team enjoys. Per-90 rates
computed over a partial season remain noisy even above the minutes threshold, which is
the motivation for the shrinkage analysis in Section~\ref{sec:results:shrinkage}. The
study covers one league and two seasons, so nothing here should be read as a claim about
football in general.
